\subsection{Đề 2}
\Opensolutionfile{ans}[ans/OTHK1-DE2-LC]
\caulc
\begin{ex}%[2D1N1-2]
	Cho hàm số $y=f(x)$ có bảng biến thiên như sau
	\begin{center}
		\begin{tikzpicture}
			\tkzTabInit[espcl=2.5,lgt=1.2]
			{$x$/0.6,$f'(x)$/0.6,$f(x)$/2}
			{$-\infty$,$-1$,0,$1$,$+\infty$}
			\tkzTabLine{,+,0,-,0,+,0,-}
			\tkzTabVar{-/$-\infty$,+/$2$,-/$1$,+/$2$,-/$-\infty$}
		\end{tikzpicture}
	\end{center}
	Hàm số đã cho đồng biến trên khoảng nào dưới đây?
	\choice
	{$(1;+\infty)$}
	{$(-1;0)$}
	{$(-1;1)$}
	{\True $(0;1)$}
	\loigiai{
		Dựa vào bảng biến thiên, ta thấy $f'(x)>0$ trên các khoảng $(-\infty;-1)$ và $(0;1)$.\\
		Do đó hàm số đồng biến trên khoảng $(0;1)$.
	}
\end{ex}

\begin{ex}%[2D1N4-1]
	Đường tiệm cận đứng của đồ thị hàm số $y=\dfrac{x+1}{x-2}$ là
	\choice
	{$y=2$}
	{\True $x=2$}
	{$y=1$}
	{$x=1$}
	\loigiai{
		Ta có $\heva{&\lim\limits_{x\to 2^{-}}f(x)=\lim\limits_{x\to 2^{-}} \dfrac{x+1}{x-2}=-\infty\\&\lim\limits_{x\to 2^{+}} f(x)=\lim\limits_{x\to 2^{+}}\dfrac{x+1}{x-2}=+\infty.}$ \\
		Vậy $x=2$ là tiệm cận đứng của đồ thị hàm số đã cho.
	}
\end{ex}

\begin{ex}%[2H2N1-1]
	Cho hình chóp $S.ABCD$ có đáy $ABCD$ là hình bình hành, tâm $O$. Khẳng định nào dưới đây \textbf{đúng}?
	\choice
	{$\overrightarrow{AC}=\overrightarrow{BD}$}
	{$\overrightarrow{AB}=\overrightarrow{DC}$}
	{\True $\overrightarrow{BC}=\overrightarrow{AD}$}
	{$\overrightarrow{OA}=\overrightarrow{OB}$}
	\loigiai{
		Do $ABCD$ là hình bình hành nên $\heva{&BC\parallel AD&\\&BC=AD\\&\overrightarrow{BC}~\text{cùng chiều với}~\overrightarrow{AD}.}$\\
		Do đó $\overrightarrow{BC}=\overrightarrow{AD}$.
	}
\end{ex}

\begin{ex}%[2H2N2-2]
	Trong không gian với hệ tọa độ $Oxyz$, cho vectơ $\overrightarrow{OM}=2\overrightarrow{i}-3\overrightarrow{k}$. Tọa độ của điểm $M$ là
	\choice
	{$M(0;2;-3)$}
	{$M(2;-3;0)$}
	{\True $M(2;0;-3)$}
	{$M(0;-3;2)$}
	\loigiai{
		Ta có $\overrightarrow{OM}=2\overrightarrow{i}-3\overrightarrow{k}=2\overrightarrow{i}+0\overrightarrow{j} -3\overrightarrow{k}\Rightarrow M(2;0;-3)$.
	}
\end{ex}

\begin{ex}%[2H2N2-3]
	Trong không gian với hệ tọa độ $Oxyz$, cho $\overrightarrow{a}=-\overrightarrow{i}+4\overrightarrow{j}+2\overrightarrow{k}$. Tọa độ của vectơ $\overrightarrow{a}$ là
	\choice
	{\True $\overrightarrow{a}=(-1;4;2)$}
	{$\overrightarrow{a}=(1;-4;-2)$}
	{$\overrightarrow{a}=(-1;-4;2)$}
	{$\overrightarrow{a}=(-1;4;-2)$}
	\loigiai{
		Ta có $\overrightarrow{a}=-\overrightarrow{i}+4\overrightarrow{j}+2\overrightarrow{k}\Rightarrow \overrightarrow{a}=(-1;4;2)$.
	}
\end{ex}

\begin{ex}%[2H2N2-2]
	Trong không gian với hệ tọa độ $Oxyz$, cho hai điểm $A(1;2;-4)$ và $B(-3;2;2)$. Tọa độ của vectơ $\overrightarrow{AB}$ là
	\choice
	{$(-2;4;-2)$}
	{\True $(-4;0;6)$}
	{$(4;0;-6)$}
	{$(-1;2;-1)$}
	\loigiai{
		Ta có $\overrightarrow{AB}=(x_B-x_A;y_B-y_A;z_B-z_A)=(-4;0;6)$. 
	}
\end{ex}

\begin{ex}%[2H2N2-2]
	Trong không gian với hệ tọa độ $Oxyz$, cho ba điểm $A(1;2;1)$, $B(2;1;3)$ và $C(0;3;2)$. Tọa độ trọng tâm $G$ của tam giác $ABC$ là
	\choice
	{$(0;6;6)$}
	{$\left(\dfrac{1}{3};\dfrac{2}{3};\dfrac{2}{3}\right) $}
	{$(3;6;6)$}
	{\True $(1;2;2)$}
	\loigiai{
		Ta có $G\left(\dfrac{x_A+x_B+x_C}{3};\dfrac{y_A+y_B+y_C}{3};\dfrac{z_A+z_B+z_C}{3}\right)=\left(1;2;2\right)$. 
	}
\end{ex}

\begin{ex}%[2D4N1-1]
	Gọi $Q_1, Q_2, Q_3$ lần lượt là giá trị của tứ phân vị thứ nhất, tứ phân vị thứ hai và tứ phân vị thứ ba của một mẫu số liệu ghép nhóm. Công thức tính khoảng tứ phân vị của mẫu số liệu đó là
	\choice
	{$\Delta_Q=2Q_2$}
	{$\Delta_Q=Q_1+Q_3$}
	{\True $\Delta_Q=Q_3-Q_1$}
	{$\Delta_Q=\dfrac{1}{3}\left( Q_1+Q_2+Q_3\right)$}
	\loigiai{
		Công thức tính khoảng tứ phân vị $\Delta_Q=Q_3-Q_1$. 
	}
\end{ex}

\begin{ex}%[2D4N1-2]
	Thời gian hoàn thành bài kiểm tra Toán $45$ phút của các bạn trong lớp được cho như sau:
	\begin{center}
		\setlength{\extrarowheight}{2pt}
		\newcommand{\PreBackslash}[1]{\let\temp=\\#1\let\\=\temp}
		\let\PBS=\PreBackslash
		\begin{tabular}{|>{\PBS\centering}m{4cm}|>{\PBS\centering}m{2cm}|>{\PBS\centering}m{2cm}| >{\PBS\centering}m{2cm}|>{\PBS\centering}m{2cm}|}\hline
			\textbf{Thời gian (phút)} &$[25;30)$&$[30;35)$&$[35;40)$&$[40;45)$\\
			\hline
			\textbf{Số học sinh}&$2$&$7$&$10$&$25$\\
			\hline
		\end{tabular}
	\end{center}
	Khoảng biến thiên về thời gian hoàn thành bài kiểm tra của các bạn trong lớp có giá trị bằng
	\choice
	{$23$}
	{$15$}
	{$10$}
	{\True $20$}
	\loigiai{
		Khoảng biến thiên $R=x_n-x_1=45-25=20$.
	}
\end{ex}

\begin{ex}%[2D4H1-3]
	Cho mẫu số liệu ghép nhóm về tuổi thọ (đơn vị tính là năm) của một loại bóng đèn mới như sau.
	\begin{center}
		\setlength{\extrarowheight}{2pt}
		\newcommand{\PreBackslash}[1]{\let\temp=\\#1\let\\=\temp}
		\let\PBS=\PreBackslash
		\begin{tabular}{|>{\PBS\centering}m{4cm}|>{\PBS\centering}m{2cm}|>{\PBS\centering}m{2cm}| >{\PBS\centering}m{2cm}|>{\PBS\centering}m{2cm}|}\hline
			\textbf{Tuổi thọ (năm)} &$[2;3{,}5)$&$[3{,}5;5)$&$[5;6{,}5)$&$[6{,}5;8)$\\
			\hline
			\textbf{Số bóng đèn}&$8$&$22$&$35$&$15$\\
			\hline
		\end{tabular}
	\end{center}
	Khoảng tứ phân vị của mẫu số liệu ghép nhóm trên bằng
	\choice
	{\True $\dfrac{303}{154}$}
	{$\dfrac{44}{7}$}
	{$\dfrac{95}{22}$}
	{$5$}
	\loigiai{
		Mẫu số liệu sắp xếp theo thứ tự không giảm: $x_1,x_2,\ldots,x_{80}$.\\
		Ta có $n=80$ nên tứ phân vị thứ nhất thuộc nhóm $[3{,}5;5)$.\\
		Công thức tính giá trị của tứ phân vị thứ $k$: $Q_k=u_m+\dfrac{\dfrac{kn}{4}-C}{n_m}\left(u_{m+1}-u_m \right)$.\\
		Do đó $Q_1=3{,}5+\dfrac{\dfrac{1\cdot 80}{4}-8}{22}\left(5-3{,}5\right)=\dfrac{95}{22}$.\\
		Ta có tứ phân vị thứ ba thuộc nhóm $[5;6{,}5)$.\\
		Do đó $Q_3=5+\dfrac{\dfrac{3\cdot 80}{4}-(8+22)}{35}\left(6{,}5-5\right)=\dfrac{44}{7}$.\\
		Vậy khoảng tứ phân vị là: $\Delta_Q=Q_3-Q_1=\dfrac{44}{7}-\dfrac{95}{22}=\dfrac{303}{154}$.
	}
\end{ex}

\begin{ex}%[2D4N1-4]
	Xét mẫu số liệu ghép nhóm cho bởi bảng sau :
	\begin{center}
		\setlength{\extrarowheight}{2pt}
		\newcommand{\PreBackslash}[1]{\let\temp=\\#1\let\\=\temp}
		\let\PBS=\PreBackslash
		\begin{tabular}{|>{\PBS\centering}m{4cm}|>{\PBS\centering}m{2cm}|>{\PBS\centering}m{2cm}| >{\PBS\centering}m{2cm}|>{\PBS\centering}m{2cm}|}\hline
			\textbf{Nhóm} &$[u_1;u_2)$&$[u_2;u_3)$&$\ldots$&$[u_k;u_{k+1})$\\
			\hline
			\textbf{Giá trị đại diện}&$c_1$&$c_2$&$\ldots$&$c_k$\\
			\hline
			\textbf{Tần số}&$n_1$&$n_2$&$\ldots$&$n_k$\\
			\hline
		\end{tabular}
	\end{center}
	Công thức tính phương sai của mẫu số liệu ghép nhóm trên là
	\choice
	{$\overline{x}= \dfrac{1}{n}\left( n_1c_1+n_2c_2+\ldots+n_kc_k\right)$}
	{\True $S^2= \dfrac{1}{n}\left( n_1c_1^2+n_2c_2^2+\ldots+n_kc_k^2\right)-\overline{x}^2$}
	{${\widehat{S}}^2= \dfrac{1}{n-1}\left[ n_1\left(c_1-\overline{x}\right) ^2+n_2\left( c_2-\overline{x}\right) ^2+\ldots+n_k\left( c_k-\overline{x}\right)^2\right]$}
	{$n=n_1+n_2+\ldots+n_k$}
	\loigiai{
		Công thức tính phương sai của mẫu số liệu ghép nhóm $$S^2= \dfrac{1}{n}\left( n_1c_1^2+n_2c_2^2+\ldots+n_kc_k^2\right)-\overline{x}^2.$$
	}
\end{ex}

\begin{ex}%[2D4H2-2]
	Bảng dưới đây thống kê lại cự li  ném tạ của một vận động viên.
	\begin{center}
		\setlength{\extrarowheight}{2pt}
		\newcommand{\PreBackslash}[1]{\let\temp=\\#1\let\\=\temp}
		\let\PBS=\PreBackslash
		\begin{tabular}{|>{\PBS\centering}m{4cm}|>{\PBS\centering}m{2cm}|>{\PBS\centering}m{2cm}| >{\PBS\centering}m{2cm}|>{\PBS\centering}m{2cm}|>{\PBS\centering}m{2cm}|}\hline
			\textbf{Cự li (mét)} &$[19;19{,}5)$&$[19{,}5;20)$&$[20;20{,}5)$&$[20{,}5;21)$&$[21;21{,}5)$\\
			\hline
			\textbf{Số lần ném}&$13$&$45$&$24$&$12$&$6$\\
			\hline
		\end{tabular}
	\end{center}
	Độ lệch chuẩn của mẫu số liệu ghép nhóm trên xấp xỉ bằng giá trị nào dưới đây?
	\choice
	{$0{,}277$}
	{$20{,}015$}
	{\True $0{,}526$}
	{$0{,}280$}
	\loigiai{
		Ta có bảng sau:
		\begin{center}
			\setlength{\extrarowheight}{2pt}
			\newcommand{\PreBackslash}[1]{\let\temp=\\#1\let\\=\temp}
			\let\PBS=\PreBackslash
			\begin{tabular}{|>{\PBS\centering}m{4cm}|>{\PBS\centering}m{2cm}|>{\PBS\centering}m{2cm}| >{\PBS\centering}m{2cm}|>{\PBS\centering}m{2cm}|>{\PBS\centering}m{2cm}|}\hline
				\textbf{Cự li (mét)} &$[19;19{,}5)$&$[19{,}5;20)$&$[20;20{,}5)$&$[20{,}5;21)$&$[21;21{,}5)$\\
				\hline
				\textbf{Giá trị đại diện}&$19{,}25$&$19{,}75$&$20{,}25$&$20{,}75$&$21{,}25$\\
				\hline
				\textbf{Số lần ném}&$13$&$45$&$24$&$12$&$6$\\
				\hline
			\end{tabular}
		\end{center}
		Cỡ mẫu $n=100$.\\
		Cự li ném trung bình của vận động viên : $$\overline{x}=\dfrac{19{,}25\cdot 13+19{,}75\cdot 45+20{,}25\cdot 24+20{,}75\cdot 12+21{,}25\cdot 6}{100}=20{,}015.$$
		Độ lệch chuẩn của mẫu số liệu ghép nhóm trên là:
		\begin{eqnarray*}
			S&=&\sqrt{S^2}\\
			&=&\sqrt{\dfrac{1}{100}\left(19{,}25^2\cdot 13+19{,}7562\cdot 45+20{,}25^2\cdot 24+20{,}75^2\cdot 12+21{,}25^2\cdot 6 \right)-20{,}015^2}\\
			&\approx&0{,}526.
		\end{eqnarray*}
	}
\end{ex}
\Closesolutionfile{ans}
\indapan{6}{ans/OTHK1-DE2-LC}
\Opensolutionfile{ans}[ans/OTHK1-DE2-DS]
\cauds
\begin{ex}%[2D1H5-7]
	Cho hàm số $y=\dfrac{x+1}{x-1}$ có đồ thị $(C)$. Khi đó
	\choiceTF
	{Đồ thị của hàm số là một đường liên tục}
	{\True Đồ thị hàm số cắt trục hoành tại điểm có hoành độ $x=-1$}
	{Hàm số nghịch biến trên $\mathbb{R}$}
	{\True Đồ thị hàm số nhận $I(1; 1)$ làm tâm đối xứng}
	\loigiai{
		Tập xác định $\mathscr{D}=\mathbb{R}\setminus \{1\}$. \\
		Đạo hàm $y'=\dfrac{-2}{(x-1)^2}<0,\, \forall x\in \mathscr{D}$.\\
		Ta có 
		\begin{itemize}
			\item $\heva{&\lim\limits_{x\to 1^+} \dfrac{x+1}{x-1} =+\infty\\ 
				&\lim\limits_{x\to 1^-} \dfrac{x+1}{x-1} =-\infty} \Rightarrow x=1$ là tiệm cận đứng.
			\item $\heva{&\lim\limits_{x\to +\infty} \dfrac{x+1}{x-1} =1\\ 
				&\lim\limits_{x\to -\infty} \dfrac{x+1}{x-1} =1} \Rightarrow y=1$ là tiệm cận ngang.
		\end{itemize}
		\begin{itemchoice}
			\itemch Do $\mathscr{D}=\mathbb{R}\setminus \{1\}$ nên hàm số không liên tục tại $x=1$.
			\itemch Ta có $\dfrac{x+1}{x-1}=0 \Rightarrow x=-1$.
			\itemch Do $\mathscr{D}=\mathbb{R}\setminus \{1\}$ nên hàm số không thể nghịch biến trên $\mathbb{R}$.
			\itemch Giao hai tiệm cận là $I(1; 1)$ nên đồ thị hàm số đối xứng qua $I$.
		\end{itemchoice}
	}
\end{ex}
\begin{ex}%[2H2H1-3]
	Cho hình lập phương $ABCD.A'B'C'D'$ cạnh $a$. Khi đó
	\choiceTF
	{\True $\left(\overrightarrow{AB}, \overrightarrow{A'D'}\right)=90^\circ$}
	{$\overrightarrow{A'C'}\cdot \overrightarrow{AD}=a^2\sqrt{2}$}
	{$\overrightarrow{AC'}= \overrightarrow{BC'}-\overrightarrow{BA'}$}
	{\True $\overrightarrow{AC'}= \overrightarrow{AB}+\overrightarrow{AA'}+\overrightarrow{AD}$}
	\loigiai{
		\begin{itemchoice}
			\itemch Ta có $\left(\overrightarrow{AB}, \overrightarrow{A'D'}\right)
			=\left(\overrightarrow{AB}, \overrightarrow{AD}\right)
			=90^\circ$.
			\itemch Ta có $\overrightarrow{A'C'}\cdot \overrightarrow{AD}
			=\overrightarrow{AC}\cdot \overrightarrow{AD}
			=a\cdot a\sqrt{2}\cdot \cos 45^\circ =a^2$.
			\itemch Ta có $\overrightarrow{AC'}= \overrightarrow{BC'}-\overrightarrow{BA}$.
			\itemch Ta có $\overrightarrow{AC'}= \overrightarrow{AB}+\overrightarrow{AA'}+\overrightarrow{AD}$.
		\end{itemchoice}
	}
\end{ex}
\begin{ex}%[2H2V2-5]
	Trong không gian với hệ trục tọa độ $Oxyz$, cho $\overrightarrow{OA}=\overrightarrow{i}-\overrightarrow{j}$, $\overrightarrow{OB}=-\overrightarrow{i}+2\overrightarrow{j}+\overrightarrow{k}$ và $C(2;2;-4)$. Khi đó
	\choiceTF
	{\True $\overrightarrow{AB}=(-2; 3; 1)$}
	{\True Trực tâm của $\triangle ABC$ có hoành độ là $\dfrac{43}{71}$}
	{Tổng các tọa độ của tâm đường tròn ngoại tiếp $\triangle ABC$ bằng $\dfrac{57}{71}$}
	{\True Cao độ đỉnh thứ tư của hình bình hành $ABCD$ là $-5$}
	\loigiai{
		Ta có $A(1; -1; 0)$, $B(-1; 2; 1)$, $C(2;2;-4)$.
		\begin{itemchoice}
			\itemch Ta có $\overrightarrow{AB}=\left(x_B-x_A; y_B-y_A; z_B-z_A\right)=(-2; 3; 1)$.
			\itemch Gọi $H(a; b; c)$ là trực tâm $\triangle ABC$. Ta có hệ phương trình
			\begin{eqnarray*}
				\heva{&\overrightarrow{AH}\cdot \overrightarrow{BC}=0 \\ &\overrightarrow{BH}\cdot \overrightarrow{AC}=0 \\ &\left[\overrightarrow{AB}, \overrightarrow{AC}\right]\cdot \overrightarrow{AH}=0}
				\Leftrightarrow \heva{&3a-5c-3=0 \\ &a+3b-4c-2=0 \\ &15a+7b+9c-8=0}
				\Leftrightarrow \heva{&a=\dfrac{43}{71} \\ &b=\dfrac{53}{355} \\ &c=-\dfrac{84}{355}.}
			\end{eqnarray*}
			Vậy trực tâm của $\triangle ABC$ có hoành độ là $\dfrac{43}{71}\cdot$
			\itemch Gọi $I(a; b; c)$ là tâm đường tròn ngoại tiếp $\triangle ABC$. Ta có
			\begin{eqnarray*}
				\heva{&AI=BI \\ &AI=CI \\ &\left[\overrightarrow{AB}, \overrightarrow{AC}\right]\cdot \overrightarrow{AI}=0}
				\Leftrightarrow \heva{&2a-3b-c=-2 \\ &a+3b-4c=11 \\ &15a+7b+9c=8}
				\Leftrightarrow \heva{&a=\dfrac{46}{71} \\ &b=\dfrac{557}{355} \\ &c=-\dfrac{501}{355}\cdot}
			\end{eqnarray*}
			Tổng các tọa độ của tâm đường tròn ngoại tiếp $\triangle ABC$ bằng $\dfrac{286}{355}\cdot$
			\itemch Ta có $ABCD$ là hình bình hành nên ta có
			\begin{eqnarray*}
				\overrightarrow{AD}=\overrightarrow{BC}
				\Leftrightarrow \heva{&x_D-1=3 \\ &y_D+1=0 \\ &z_D-0=-5}
				\Leftrightarrow \heva{&x_D=4 \\ &y_D=-1 \\ &z_D=-5.}
			\end{eqnarray*}
			Vậy cao độ của đỉnh thứ tư hình bình hành $ABCD$ bằng $-5$.
		\end{itemchoice}
	}
\end{ex}
\begin{ex}%[2D3B2-2]
	Cho mẫu số liệu ghép nhóm cho bởi bảng sau:
	\begin{center}
		\begin{tabular}{|c|c|c|c|c|}
			\hline Nhóm & $\left[u_1; u_2\right)$ & $\left[u_2; u_3\right)$ & $\ldots$ & $\left[u_k; u_{k+1}\right)$ \\
			\hline Giá trị đại diện& $c_1$ & $c_2$ & $\ldots$ & $c_k$ \\
			\hline Tần số & $n_1$ & $n_2$ & $\ldots$ & $n_k$ \\
			\hline
		\end{tabular}
	\end{center}
	\choiceTF
	{Số trung bình đại diện cho nửa mẫu số liệu dưới}
	{\True Phương sai $S^2=\dfrac{1}{n}\sum\limits_{i=1}^k\, \dfrac{1}{n_i}\left(c_i-\overline{x}\right)^2$ với $n=\sum\limits_{i=1}^k\, n_i$}
	{\True Phương sai dùng để đo mức độ phân tán của mẫu số liệu ghép nhóm quanh giá trị trung bình}
	{Độ lệch chuẩn càng lớn thì mức độ phân tán của mẫu số liệu càng nhỏ}
	\loigiai{
		\begin{itemchoice}
			\itemch Số trung bình đại diện cho cả mẫu số liệu.
			\itemch Phương sai $S^2=\dfrac{1}{n}\sum\limits_{i=1}^k\, \dfrac{1}{n_i}\left(c_i-\overline{x}\right)^2$ với $n=\sum\limits_{i=1}^k\, n_i$.
			\itemch Phương sai (độ lệch chuẩn) dùng để đo mức độ phân tán của mẫu số liệu ghép nhóm quanh giá trị trung bình.
			\itemch Độ lệch chuẩn càng lớn thì mức độ phân tán của mẫu số liệu càng lớn.
		\end{itemchoice}
	}
\end{ex}
\Closesolutionfile{ans}
\indapan{3}{ans/OTHK1-DE2-DS}
\Opensolutionfile{ans}[ans/OTHK1-DE2-KQ]
\caukq
\begin{ex}%[2D1H5-1]
	Đường cong ở hình bên dưới là đồ thị của hàm số có dạng $y=\dfrac{-x+b}{cx+d}$.
	\begin{center}
		\begin{tikzpicture}[scale=0.8,font=\footnotesize,line join=round, line cap=round,>=stealth]
			\draw[->] (-5.1,0)--(4.1,0) node[below left] {$x$};
			\draw[->] (0,-5.1)--(0,4.1) node[below left] {$y$};
			\draw (0,0) node [below left] {$O$};
			\foreach \x/\nx in {-1/-1,1/1}
			\draw[thin] (\x,1pt)--(\x,-1pt) node [below left] {$\nx$};
			\foreach \y/\ny in {-1/-1,1/1}
			\draw[thin] (1pt,\y)--(-1pt,\y) node [below left] {$\ny$};
			\draw[dashed,thin] (-0.99,-5)--(-0.99,4);
			\begin{scope}
				\clip (-5,-5) rectangle (4,4);
				\draw[thick,samples=200,domain=-5:-1.01,smooth,variable=\x] plot (\x,{(-1*(\x)+1)/(1*(\x)+1)});
				\draw[thick,samples=200,domain=-0.99:4,smooth,variable=\x] plot (\x,{(-1*(\x)+1)/(1*(\x)+1)});
				\draw[dashed,thin] (-5,-1/1)--(4,-1/1);
			\end{scope}
		\end{tikzpicture}
	\end{center}
	Tính giá trị của biểu thức $P=b+c+d$.
	\shortans[]{$3$}
	\loigiai{
		Dựa vào đồ thị, dễ thấy đồ thị có đường tiệm cận đứng $x=-1$ và đường tiệm cận ngang $y=-1$.\\
		Ta có $\heva{&\dfrac{-1}{c}=-1\\&-\dfrac{d}{c}=-1}\Rightarrow \heva{&c=1\\&d=1}$.\\
		Đồ thị đi qua điểm $(0;1)\Rightarrow b=d=1$.\\
		Vậy $P=b+c+d=1+1+1=3$.
	}
\end{ex}

\begin{ex}%[2H2N2-3]
	Trong không gian với hệ tọa độ $Oxyz$, cho điểm $A(7;-1;2)$. Vectơ $\overrightarrow{OA}$ được biểu diễn dưới dạng $\overrightarrow{OA}=a\overrightarrow{i}+b\overrightarrow{j}+c\overrightarrow{k}$. Tính giá trị của biểu thức $P=a-b+c$.
	\shortans[]{$10$}
	\loigiai{
		Ta có $A(7;-1;2)\Leftrightarrow \overrightarrow{OA}=(7;-1;2)=7\overrightarrow{i}-\overrightarrow{j}+2\overrightarrow{k}$.\\
		Vậy $a=7$;$b=-1$;$c=2\Rightarrow P=a-b+c=7-(-1)+2=10$.
	}
\end{ex}

\begin{ex}%[2D4V1-3]
	Cân nặng của một số lợn con mới sinh thuộc hai giống $A$ và $B$ được cho ở biểu đồ dưới đây (đơn vị: $\mathrm{kg}$).
	\begin{center}
		\begin{tikzpicture}[scale=1.2, font=\footnotesize, line join=round, line cap=round, >=stealth]
			\coordinate [label= left: $0$](O) at (0,0) ;
			\coordinate [label= left: (Số con)](y) at (0,6) ;
			\coordinate [label= below: (kg)](x) at (8,0) ;
			\foreach \i in {1,2,3,4,5}
			{\draw[-,color=gray!40] (0,\i)--(8,\i); }
			\draw (0,1)node[left]{$10$}(0,2)node[left]{$20$}(0,3)node[left]{$30$}(0,4)node[left]{$40$}(0,5)node[left]{$50$};
			%--------------------------------------------
			\foreach \i/\j/\a in 
			{0.4/0.8/8,1.2/1.3/13,2/2.8/28,2.8/1.4/14,3.6/3.2/32,4.4/2.4/24,	5.2/1.7/17,6/1.4/14}{
				\draw[] ($(\i,\j)+(0,0.2)$) node[scale=0.8]{$\a$};
			}
			%-----------------------------------------------
			\foreach \i/\j in 
			{0/0.8,1.6/2.8,3.2/3.2,4.8/1.7}{
				\fill[blue!50] (\i,0) rectangle+ (0.8,\j);
				\draw[thick] (\i,0) rectangle+ (0.8,\j);
			}
			%-----------------------------------------------	
			\foreach \i/\j in 
			{0.8/1.3,2.4/1.4,4/2.4,5.6/1.4}{
				\fill[pattern=north east lines] (\i,0) rectangle+ (0.8,\j);
				\draw[thick] (\i,0) rectangle+ (0.8,\j);
			}	
			%-----------------------------------------------
			\foreach \i/\j in 
			{0.4/[1{,}0;1{,}1),2/[1{,}1;1{,}2),3.6/[1{,}2;1{,}3),5.2/[1{,}3;1{,}4)}{
				\draw($(\i,0)+(0.5,0)$) node[below]{\j};
			}	
			\foreach \i/\j in {}{\draw (\i,\j) node[above]{\j};
			}			
			\draw[->,very thick] (O)--(x);
			\draw[->,very thick] (O)--(y);
			\fill[blue!50,draw] (5.5,4.2) rectangle (6,4.8);
			\fill[pattern=north east lines,draw] (5.5,3.2) rectangle (6,3.8);
			\draw (6,4.5) node [right]{Giống $A$};
			\draw (6,3.5) node [right]{Giống $B$};
			\draw (4,5.5) node [above]{\textbf{Cân nặng của một số lợn con mới sinh}};
		\end{tikzpicture}
	\end{center}
	Gọi $\Delta Q_1$, $\Delta Q_2$ lần lượt là giá trị khoảng tứ phân vị của lợn con giống $A$ và $B$. Tìm giá trị của $\Delta Q_2-\Delta Q_1$ (làm tròn kết quả đến hàng phần trăm). 
	\shortans[]{$0{,}03$}
	\loigiai{
		Cân nặng của lợn con giống $A$ và giống $B$ được thống kê như sau:
		\begin{center}
			\setlength{\extrarowheight}{2pt}
			\newcommand{\PreBackslash}[1]{\let\temp=\\#1\let\\=\temp}
			\let\PBS=\PreBackslash
			\begin{tabular}{|>{\PBS\centering}m{4cm}|>{\PBS\centering}m{2cm}|>{\PBS\centering}m{2cm}| >{\PBS\centering}m{2cm}|>{\PBS\centering}m{2cm}|}\hline
				\textbf{Cân nặng ($\mathrm{kg}$)} &$[1{,}0;1{,}1)$&$[1{,}1;1{,}2)$&$[1{,}2;1{,}3)$&$[1{,}3;1{,}4)$\\
				\hline
				\textbf{Giá trị đại diện}&$1{,}05$&$1{,}15$&$1{,}25$&$1{,}35$\\
				\hline
				\textbf{Số con giống $A$}&$8$&$28$&$32$&$17$\\
				\hline
				\textbf{Số con giống $B$}&$13$&$14$&$24$&$14$\\
				\hline
			\end{tabular}
		\end{center} 
		\begin{enumerate}[$\bullet$]
			\item Xét mẫu số liệu của giống lợn $A$.\\
			Ta có $n=85$ nên tứ phân vị thứ nhất thuộc nhóm $[1{,}1;1{,}2)$.\\
			Do đó $Q_1=1{,}1+\dfrac{\dfrac{85}{4}-8}{28}\cdot (1{,}2-1{,}1)=\dfrac{257}{224}$.\\
			Tứ phân vị thứ ba thuộc nhóm $[1{,}2;1{,}3)$.\\
			Do đó $Q_3=1{,}2+\dfrac{\dfrac{3\cdot 85}{4}-(8+28)}{32}\cdot (1{,}3-1{,}2)=\dfrac{1\,647}{1\,280}$.\\
			Vậy $\Delta Q_1=Q_3-Q_1=\dfrac{1\,249}{8\,960}.$
			\item Xét mẫu số liệu của giống lợn $B$.\\
			Ta có $n=65$ nên tứ phân vị thứ nhất thuộc nhóm $[1{,}1;1{,}2)$.\\
			Do đó $Q_1=1{,}1+\dfrac{\dfrac{65}{4}-13}{14}\cdot (1{,}2-1{,}1)=\dfrac{629}{560}$.\\
			Tứ phân vị thứ ba thuộc nhóm $[1{,}2;1{,}3)$.\\
			Do đó $Q_3=1{,}2+\dfrac{\dfrac{3\cdot 65}{4}-(13+14)}{24}\cdot (1{,}3-1{,}2)=\dfrac{413}{320}$.\\
			Vậy $\Delta Q_2=Q_3-Q_1=\dfrac{75}{448}$.\\
			Giá trị của $\Delta Q_2-\Delta Q_1\approx 0,03$.
		\end{enumerate}
	}
\end{ex}

\begin{ex}%[2D4V2-2]
	Thâm niên công tác của các công nhân hai nhà máy $A$ và $B$ được thống kê ở bảng dưới đây:
	\begin{center}
		\setlength{\extrarowheight}{2pt}
		\newcommand{\PreBackslash}[1]{\let\temp=\\#1\let\\=\temp}
		\let\PBS=\PreBackslash
		\begin{tabular}{|>{\PBS\centering}m{6cm}|>{\PBS\centering}m{1.5cm}|>{\PBS\centering}m{1.5cm}| >{\PBS\centering}m{1.5cm}|>{\PBS\centering}m{1.5cm}|>{\PBS\centering}m{1.5cm}|}\hline
			\textbf{Thâm niên công tác (năm)} &$[0;5)$&$[5;10)$&$[10;15)$&$[15;20)$&$[20;25)$\\
			\hline
			\textbf{Số công nhân nhà máy $A$}&$35$&$13$&$12$&$12$&$8$\\
			\hline
			\textbf{Số công nhân nhà máy $B$}&$14$&$26$&$24$&$11$&$5$\\
			\hline
		\end{tabular}
	\end{center}
	Gọi $S_A$, $S_B$ lần lượt là độ lệch chuẩn về thâm niên công tác của công nhân hai nhà máy $A$ và $B$. Tính giá trị của $S_A-S_B$ (làm tròn kết quả đến hàng phần trăm).
	\shortans[]{$1{,}48$}
	\loigiai{
		Thâm niên công tác của công nhân hai nhà máy $A$ và $B$ được thống kê như sau:
		\begin{center}
			\setlength{\extrarowheight}{2pt}
			\newcommand{\PreBackslash}[1]{\let\temp=\\#1\let\\=\temp}
			\let\PBS=\PreBackslash
			\begin{tabular}{|>{\PBS\centering}m{6cm}|>{\PBS\centering}m{1.5cm}|>{\PBS\centering}m{1.5cm}| >{\PBS\centering}m{1.5cm}|>{\PBS\centering}m{1.5cm}|>{\PBS\centering}m{1.5cm}|}
				\hline
				\textbf{Thâm niên công tác (năm)} &$[0;5)$&$[5;10)$&$[10;15)$&$[15;20)$&$[20;25)$\\
				\hline
				\textbf{Giá trị đại diện} &$2{,}5$&$7{,}5$&$12{,}5$&$17{,}5$&$22{,}5$\\
				\hline
				\textbf{Số công nhân nhà máy $A$}&$35$&$13$&$12$&$12$&$8$\\
				\hline
				\textbf{Số công nhân nhà máy $B$}&$14$&$26$&$24$&$11$&$5$\\
				\hline
			\end{tabular}
		\end{center}
		\begin{enumerate}[$\bullet$]
			\item Xét mẫu số liệu thâm niên công tác của công nhân nhà máy $A$.\\
			Thâm niên trung bình của công nhân nhà máy $A$ :
			$$\overline{x}_A=\dfrac{1}{80}\left(2{,}5\cdot 35+7{,}5\cdot 13+12{,}5\cdot 12+17{,}5\cdot 12+22{,}5\cdot8\right)=\dfrac{145}{16}.$$
			Phương sai về thâm niên của công nhân nhà máy $A$: $$S_A^2=\dfrac{1}{80}\left(2{,}5^2\cdot 35+7{,}5^2\cdot 13+12{,}5^2\cdot 12+17{,}5^2\cdot 12+22{,}5^2\cdot8\right)-\left(\dfrac{145}{16}\right)^2 =\dfrac{12\,735}{256}.$$
			Độ lệch chuẩn  về thâm niên của công nhân nhà máy $A$: $S_A=\sqrt{S_A^2}\approx 7{,}05$.
			\item Xét mẫu số liệu thâm niên công tác của công nhân nhà máy $B$.\\
			Thâm niên trung bình của công nhân nhà máy $B$ :
			$$\overline{x}_B=\dfrac{1}{80}\left(2{,}5\cdot 14+7{,}5\cdot 26+12{,}5\cdot 24+17{,}5\cdot 11+22{,}5\cdot5\right)=\dfrac{167}{16}.$$
			Phương sai về thâm niên của công nhân nhà máy $B$: $$S_B^2=\dfrac{1}{80}\left(2{,}5^2\cdot 14+7{,}5^2\cdot 26+12{,}5^2\cdot 24+17{,}5^2\cdot 11+22{,}5^2\cdot 5\right)-\left(\dfrac{167}{16}\right)^2 =\dfrac{7\,951}{256}.$$
			Độ lệch chuẩn  về thâm niên của công nhân nhà máy $B$: $S_B=\sqrt{S_B^2}\approx 5{,}57$.\\
			Giá trị của $S_A-S_B\approx 7{,}05-5{,}57\approx 1{,}48$.
		\end{enumerate}
	}
\end{ex}

\begin{ex}%[2H2C1-2]
	Trong không gian $Oxyz$, cho hình chóp $S.ABCD$ có đáy $ABCD$ là hình thoi cạnh bằng $a$, giao điểm hai đường chéo $AC$ và $BD$ trùng với gốc $O$. Các vectơ $\overrightarrow{OB}$, $\overrightarrow{OC}$, $\overrightarrow{OS}$ lần lượt cùng hướng với $\overrightarrow{i}$, $\overrightarrow{j}$, $\overrightarrow{k}$ và $OA=OS=4$. Giả sử $\overrightarrow{AB}=x_1\overrightarrow{i}+y_1\overrightarrow{j}+z_1\overrightarrow{k}$; $\overrightarrow{AC}=x_2\overrightarrow{i}+y_2\overrightarrow{j}+z_2\overrightarrow{k}$ và $\overrightarrow{SB}=x_3\overrightarrow{i}+y_3\overrightarrow{j}+z_3\overrightarrow{k}$. Tính giá trị của $a$ để biểu thức $P=x_1+x_2+x_3+y_1+y_2+y_3+z_1+z_2+z_3$ đạt giá trị nhỏ nhất.
	\shortans[]{$4$}
	\loigiai{
		\begin{center}
			\begin{tikzpicture}[scale=1,font=\footnotesize,line join=round, line cap=round,>=stealth]
				\coordinate (A) at (0,0);
				\coordinate (B) at (-2,-2);
				\coordinate (D) at (5,0);
				\coordinate (C) at ($(B)+(D)-(A)$);
				\coordinate (O) at ($(A)!0.5!(C)$);
				\coordinate (S) at ($(O)+(0,5)$);
				\draw[thick] (S)--(B) (S)--(C)  (B)--(C)(S)--(D)(C)--(D);
				\draw[dashed,thin](S)--(A) (A)--(B) (A)--(C) (A)--(D)   (S)--(O) (B)--(D);
				\pic[draw,thin,angle radius=2mm] {right angle = S--O--D} pic[draw,thin,angle radius=2mm] {right angle = S--O--A}pic[draw,thin,angle radius=2mm] {right angle = C--O--D};
				\draw[dashed,->] (O)--($(O)+(0,0.9)$)node [right]{$\overrightarrow{k}$};
				\draw[dashed,->] (O)--($(O)+(-1,-0.28)$)node [below]{$\overrightarrow{i}$};
				\draw[dashed,->] (O)--($(O)+(0.8,-0.53)$)node [above]{$\overrightarrow{j}$};
				\foreach \i/\g in {S/90,A/180,B/-90,C/-90,D/0,O/-90}{\draw[fill=white](\i) circle (0pt) ($(\i)+(\g:3mm)$) node[scale=1]{$\i$};}
			\end{tikzpicture}
		\end{center}
		Ta có $OA=OC=OS=4\Rightarrow \heva{&\overrightarrow{OA}=-4\overrightarrow{j}\\&\overrightarrow{OC}=4\overrightarrow{j}\\&\overrightarrow{OS}=4\overrightarrow{k}.}$\\
		Lại có $OB=\sqrt{BC^2-OC^2}=\sqrt{a^2-4^2}=\sqrt{a^2-16}\Rightarrow \overrightarrow{OB}=\sqrt{a^2-16}\cdot\overrightarrow{i}$.\\
		Ta có 
		\begin{enumerate}[$\bullet$]
			\item $\overrightarrow{AB}=\overrightarrow{OB}-\overrightarrow{OA}=\sqrt{a^2-16}\cdot\overrightarrow{i}+4\overrightarrow{j}\Rightarrow \heva{&x_1=\sqrt{a^2-16}\\&y_1=4\\&z_1=0.}$
			\item $\overrightarrow{AC}=\overrightarrow{OC}-\overrightarrow{OA}=4\overrightarrow{j}+4\overrightarrow{j}=8\overrightarrow{j}\Rightarrow \heva{&x_2=0\\&y_2=8\\&z_2=0.}$
			\item $\overrightarrow{SB}=\overrightarrow{OB}-\overrightarrow{OS}=\sqrt{a^2-16}\cdot\overrightarrow{i}-4\overrightarrow{k}\Rightarrow \heva{&x_3=\sqrt{a^2-16}\\&y_3=0\\&z_3=-4.}$
		\end{enumerate}
		Khi đó $P=x_1+x_2+x_3+y_1+y_2+y_3+z_1+z_2+z_3=8+2\sqrt{a^2-16}$.\\
		Xét hàm số $P=8+2\sqrt{a^2-16}$.\\
		Tập xác đinh $\mathscr{D}=[4;+\infty)$ (do $a>0$).\\
		Ta có $P'=\dfrac{2a}{\sqrt{a^2-16}}>0,\,\forall a \in [4;+\infty)$.\\
		Vậy $\min\limits_{[4;+\infty)}P=P(4)=8$. 
	}
\end{ex}
\begin{ex}%[2H2C2-4]
	Trong không gian với hệ tọa độ $Oxyz$, cho vectơ $\overrightarrow{a}=(1;-2;4)$ và $\overrightarrow{b}=(x_0;y_0;z_0)$ cùng phương với $\overrightarrow{a}$. Biết $\overrightarrow{b}$ tạo với tia $Oy$ một góc nhọn và $\left| \overrightarrow{b}\right|=\sqrt{21}$. Tính giá trị của biểu thức  $P=x_0+y_0+z_0$. 
	\shortans[]{$-3$}
	\loigiai{
		Vì $\overrightarrow{b}$ cùng phương với $\overrightarrow{a}\Rightarrow \overrightarrow{b}=k\overrightarrow{a}\Rightarrow \heva{&x_0=k\\&y_0=-2k\\&z_0=4k}$.\\
		Lại có $\left| \overrightarrow{b}\right|=\sqrt{21}\Leftrightarrow \sqrt{k^2+(-2k)^2+(4k)^2}=\sqrt{21}\Leftrightarrow k=\pm 1$.
		\begin{enumerate}[$\bullet$]
			\item Với $k=1\Rightarrow \overrightarrow{b}=(1;-2;4)$.\\
			Khi đó $\cos\left(\overrightarrow{b},Oy\right)=\dfrac{\overrightarrow{b}\cdot \overrightarrow{j}}{\left|\overrightarrow{b}\right|\cdot \left|\overrightarrow{j}\right|}$, trong đó $\overrightarrow{b}\cdot \overrightarrow{j}=1\cdot 0+ (-2)\cdot 1+4\cdot 0=-2<0$.\\
			Suy ra góc tạo bởi $\overrightarrow{b}$ và $Oy$ là góc tù.
			\item Với $k=-1\Rightarrow \overrightarrow{b}=(-1;2;-4)$.\\
			Khi đó $\cos\left(\overrightarrow{b},Oy\right)=\dfrac{\overrightarrow{b}\cdot \overrightarrow{j}}{\left|\overrightarrow{b}\right|\cdot \left|\overrightarrow{j}\right|}$, trong đó $\overrightarrow{b}\cdot \overrightarrow{j}=-1\cdot 0+2\cdot 1+(-4)\cdot 0=2>0$.\\
			Suy ra góc tạo bởi $\overrightarrow{b}$ và $Oy$ là góc nhọn.\\
			Vậy $k=-1$ thỏa mãn.\\
			Do đó $\overrightarrow{b}=(-1;2;-4)\Rightarrow x_0=-1$, $y_0=2$, $z_0=-2$.\\
			Vậy $P=x_0+y_0+z_0=-1+2+(-4)=-3$.
		\end{enumerate}
	}
\end{ex}
\Closesolutionfile{ans}
\indapan{6}{ans/OTHK1-DE2-KQ}